\section{Introduction}
\label{sec:intro}

Cooperative multi-agent reinforcement learning (c-MARL) trains teams of agents
that share a common reward, and is increasingly considered for systems such as
robot swarms, autonomous vehicles, and sensor networks. Most methods assume that
every teammate executes the learned cooperative policy. In practice an agent may
be compromised or malfunction and behave arbitrarily, a \emph{Byzantine} failure
\citep{lamport1982byzantine}. Because cooperative policies are trained to rely on
their teammates, even a single such agent can collapse team performance.

A principled response is to model this as a Bayesian game. Each agent is secretly
assigned a type, cooperative or adversarial. Cooperative agents do not know who is
compromised, so they maintain a \emph{belief} over their teammates' types, update
it from observed behavior, and act on it. The resulting solution concept, an
\emph{ex interim} robust Markov perfect Bayesian equilibrium, lets agents defend
only when the evidence warrants it, rather than treating every teammate as a
potential threat \citep{li2024byzantine,kazari2026bayesian}. Existing belief-based
methods, however, assume that \emph{at most one} agent is compromised. Methods that
tolerate several adversaries do so \emph{ex ante}, without inference over who they
are \citep{phan2021radar}. Belief-based robustness against multiple adversaries
remains open.

A second difficulty is computational. The robust equilibrium is a saddle point:
the team maximizes the shared return while the adversary minimizes it. Current
methods find it with two-timescale gradient descent--ascent on policy-gradient
objectives. In min--max problems such dynamics are known to cycle rather than
converge \citep{mertikopoulos2018cycles}, and we observe exactly this behavior: the
team's reward under attack oscillates and the final result depends on the seed and
on when training stops. These methods also offer no observable signal that an
equilibrium has been reached. Reward measured against a co-trained adversary can
overstate robustness when that adversary is itself stale or cycling.

\paragraph{Objective.}
This work aims to make belief-based robust c-MARL both \emph{stable to train} and
\emph{applicable to multiple adversaries}. Our approach is to replace saddle-point
dynamics with the minimization of a single non-negative \emph{energy} that vanishes
exactly at equilibrium. We build this energy from per-state optimality conditions.
For each player, the Karush--Kuhn--Tucker conditions of choosing an action
distribution reduce to sign and complementarity conditions on the advantage
function. Under entropy regularization they reduce further to the residual
$Q(s,a) - \tau \log \pi(a \mid s) - V(s) = 0$, whose joint solution is a quantal
response equilibrium that approaches the robust equilibrium as $\tau \to 0$
\citep{mckelvey1995quantal,nachum2017pcl}. The energy aggregates these residuals for
the team and the adversary. Cooperative agents' residuals are evaluated on
belief-averaged advantages, so the conditions reflect what an agent can actually
know (ex interim rather than oracle optimality). A Bellman-consistency term keeps
the critic faithful. Because training is now a single minimization, there are no
rotational dynamics to cycle, and the energy's value serves as a convergence
certificate.

With a stable training procedure in hand, we extend the setting to $k$ independent,
non-colluding adversaries among $N$ agents, using a factored belief that scales
linearly in $N$. This lets us ask how inference and robustness degrade as the
fraction $k/N$ of compromised agents grows, and whether learned teams exhibit a
tolerance threshold analogous to the classical Byzantine fault-tolerance bound.

\paragraph{Contributions.}
% TODO: tighten once results exist; keep claims matched to what is shown.
\begin{itemize}
  \item A Bayesian formulation of robust c-MARL with $k$ independent Byzantine
        agents, and the corresponding ex interim robust equilibrium.
  \item An energy-based training method built from per-state (soft) KKT residuals,
        belief-averaged advantages, and Bellman consistency, which provides a
        convergence certificate.
  \item A factored belief architecture whose cost is linear in the number of agents.
  \item An empirical study comparing the method with gradient descent--ascent and
        extragradient baselines under evaluation against freshly trained
        adversaries, and characterizing robustness as a function of $k/N$.
\end{itemize}
